%% modular_shadow_LMP.tex
%% Standalone note for Letters in Mathematical Physics
%% "The Modular Shadow Conjecture: MSS Chaos Bound from Type-II Crossed Products"
%% Extracted / elaborated from ECI v6.0.44 §sec:mss-shadow
%% Status: DRAFT 2026-05-04
%% All arXiv IDs verified against the ECI v6.0.44 bibliography.
%% The conjecture is explicitly flagged [CONJECTURE] throughout.

\documentclass[12pt,a4paper]{article}

\usepackage{amsmath,amssymb,amsthm}
\usepackage{hyperref}
\usepackage{cite}
\usepackage{geometry}
\geometry{margin=2.5cm}

%% Theorem environments
\newtheorem{conjecture}{Conjecture}
\newtheorem{remark}{Remark}
\newtheorem{proposition}{Proposition}

\title{\textbf{The Modular Shadow Conjecture}:\\[4pt]
       MSS Chaos-Bound Saturation as a Signature\\
       of Type-II$_\infty$ Crossed-Product Structure}

\author{%
  \textit{(Author names to be determined)}\\[2pt]
  \small \textit{Submitted to Letters in Mathematical Physics}
}

\date{Preprint: \today}

\begin{document}

\maketitle

\begin{abstract}
We propose the \emph{Modular Shadow Conjecture}: in the type-II$_\infty$
crossed-product von Neumann algebra associated with an observer's causal
diamond at KMS inverse temperature $\beta = 2\pi$, the Krylov--Lanczos
complexity growth saturates the Maldacena--Shenker--Stanford (MSS)
chaos bound, yielding $\lambda_L^{\mathrm{Krylov}} = 2\pi/\beta = 1$
in natural units.
We argue that this saturation is \emph{not} a generic feature of
strongly-chaotic quantum systems but is a ``modular shadow'': a
structural fingerprint of the underlying type-II$_\infty$ algebraic
data, with the KMS condition at $\beta=2\pi$ (Bisognano--Wichmann)
providing the natural temperature precisely equal to the MSS bound.
We support the conjecture with two tractable models (free scalar field
in the Rindler wedge and the Sachdev--Ye--Kitaev model at large $N$)
and discuss implications for algebra-type-dependent saturation and
the Connes--Stoermer trace on the type-II$_1$ reduction.
[CONJECTURE: not yet proven as a theorem.]
Throughout this paper, $\rho$ denotes the dimensionless analog-Hawking saturation ratio
(as in the companion Cardy paper~\cite{CardyRho}); density-matrix notation is not used.
\end{abstract}

% ─────────────────────────────────────────────────────────────────────
\section{Background}
\label{sec:background}
% ─────────────────────────────────────────────────────────────────────

\paragraph{MSS chaos bound.}
Maldacena, Shenker, and Stanford~\cite{MSS2016} proved that in any
quantum system with a large-$N$ limit and a thermal density matrix at
inverse temperature $\beta$, the out-of-time-order correlator (OTOC)
Lyapunov exponent satisfies
\begin{equation}
  \lambda_L \;\le\; \frac{2\pi k_B T}{\hbar}
            \;=\; \frac{2\pi}{\beta}
  \label{eq:mss}
\end{equation}
with equality in maximally chaotic systems such as the
Sachdev--Ye--Kitaev (SYK) model and near-extremal black holes.
The bound~\eqref{eq:mss} is usually interpreted as a kinematic
consequence of analyticity and unitarity in the large-$N$ limit,
without reference to the algebraic structure of the observable algebra.

\paragraph{Krylov complexity.}
Parker, Cao, Avdoshkin, Scaffidi, and Altman~\cite{Parker2019}
introduced Krylov (Lanczos) complexity $\mathcal{C}_K$ as a
basis-independent measure of operator growth: given an operator
$\mathcal{O}$ evolved by a Hamiltonian $H$, one expands
$e^{iHt}\mathcal{O}e^{-iHt}$ in a basis of orthogonal
``Krylov--Koenig'' vectors, with complexity equal to the mean
position in Krylov space.
Parker et al.\ conjectured and tested that in thermal systems
$\mathcal{C}_K(t) \sim e^{\lambda_K t}$ at late times, and found
numerically that $\lambda_K \to 2\pi/\beta$ in maximally chaotic
models, saturating the MSS bound~\eqref{eq:mss}.
Caputa, Mag\'{a}n, Patramanis, and Tonni~\cite{Caputa2024} subsequently
showed that a \emph{modular} version of Krylov complexity, constructed
from the modular Hamiltonian rather than from the physical Hamiltonian,
exhibits a universal Lyapunov exponent $\lambda_L^{\mathrm{mod}} = 2\pi$
in dimensionless modular time for conformal field theories (CFTs),
providing the first algebraic bridge between Krylov growth and
modular evolution.

\paragraph{Type-II$_\infty$ crossed products.}
Witten~\cite{Witten2022} and Chandrasekaran--Longo--Penington--Witten~\cite{CLPW2023}
showed that the algebra of observables of a gravitational system in the
presence of a quantum reference frame (QRF) is promoted from a
type-III$_1$ von Neumann algebra (the standard Rindler or wedge algebra)
to a type-II$_\infty$ crossed-product algebra:
\begin{equation}
  \mathcal{A}_{\mathrm{II}_\infty}
  \;=\; \mathcal{A}_{\mathrm{III}_1} \rtimes_{\sigma^\phi} \mathbb{R},
  \label{eq:crossed}
\end{equation}
where $\sigma^\phi_t$ is the modular automorphism group of the KMS
state $\phi$.
Crucially, the type-II$_\infty$ algebra \emph{admits a semifinite
trace} $\tau$, which is absent in the type-III$_1$ case; this is what
allows well-defined notions of generalised entropy and observer-dependent
gravitational entropy~\cite{DEHK2025a,DEHK2025b}.

De Vuyst, Eccles, H\"{o}hn, and Kirklin~\cite{DEHK2025a} established that
the observer-dependent gravitational entropy in the crossed-product
setting naturally yields a finite von Neumann entropy on the
QRF-dressed algebra, bypassing the type-III ultraviolet divergences of
the bare algebra.
Faulkner and Speranza~\cite{FaulknerSperanza2024} connected the modular
flow on $\mathcal{A}_{\mathrm{II}_\infty}$ to the generalised second law
(GSL), showing that the modular Hamiltonian $K_\phi$ generates the
evolution that computes the generalised entropy.

\paragraph{KMS condition and $\beta = 2\pi$.}
The Bisognano--Wichmann theorem (1976) establishes that the vacuum
state of a relativistic QFT, restricted to the Rindler wedge
$\mathcal{W} = \{x^1 > |x^0|\}$, is a KMS state at inverse temperature
$\beta_{\mathrm{BW}} = 2\pi$ with respect to the boost generator
$K_{\mathrm{boost}}$~\cite{BisognanoWichmann1976}.
In the observer's static diamond (or Rindler wedge), the modular
Hamiltonian is proportional to the boost:
$K_\phi = 2\pi\,K_{\mathrm{boost}}$,
and the associated surface gravity is $\kappa_R = 2\pi T_H$
(in units $k_B = \hbar = 1$).
Thus $\beta = 2\pi$ is the \emph{canonical} KMS temperature for
the causal-wedge algebra, and is singled out by the algebraic
structure rather than by any dynamical assumption.

% ─────────────────────────────────────────────────────────────────────
\section{The Modular Shadow Conjecture}
\label{sec:conjecture}
% ─────────────────────────────────────────────────────────────────────

We now state the central conjecture of this note.
Let $(\mathcal{A},\phi)$ be the type-II$_\infty$ crossed-product
algebra of an observer's causal diamond, equipped with the KMS state
$\phi$ at $\beta = 2\pi$ given by the Bisognano--Wichmann construction.
Let $K_\phi \in \mathcal{A}$ be the modular Hamiltonian, and let
$\mathcal{C}_K^{\mathrm{mod}}(s)$ denote the Krylov complexity of an
operator evolved by the modular flow:
\begin{equation}
  \mathcal{O}(s) \;=\; e^{iK_\phi s}\,\mathcal{O}\,e^{-iK_\phi s},
  \quad s \in \mathbb{R}.
  \label{eq:modular-flow}
\end{equation}

\begin{conjecture}[Modular Shadow Conjecture]\label{conj:msc}
\emph{[CONJECTURE]}
In the type-II$_\infty$ crossed-product algebra
$\mathcal{A}_{\mathrm{II}_\infty} = \mathcal{A}_{\mathrm{III}_1}
\rtimes_{\sigma^\phi} \mathbb{R}$ at KMS temperature $\beta = 2\pi$,
the modular Krylov complexity $\mathcal{C}_K^{\mathrm{mod}}(s)$
grows exponentially at late modular times with Lyapunov exponent
\begin{equation}
  \lambda_L^{\mathrm{Krylov}}
  \;=\; \frac{2\pi}{\beta}
  \;=\; 1,
  \label{eq:conj-main}
\end{equation}
thereby saturating the MSS chaos bound~\eqref{eq:mss} exactly.
Moreover, this saturation is a consequence of the type-II$_\infty$
algebraic structure and the KMS condition, not of the particular
Hamiltonian or interaction: it holds for any local algebra in the
crossed-product class~\eqref{eq:crossed}.
\end{conjecture}

\begin{remark}
The identification $\lambda_L^{\mathrm{Krylov}} = 2\pi/\beta$
uses the natural units for modular flow in which $s$ is dimensionless
modular time. In physical units with inverse temperature $\beta$,
the Krylov Lyapunov exponent $\lambda_L = 2\pi/\beta$ matches the
MSS bound exactly.
\end{remark}

\begin{remark}
We call this a ``modular shadow'' because the $2\pi T_H$ prefactor
in both the MSS bound and the Krylov saturation is a \emph{shadow}
of the underlying modular structure: it is the type-II$_\infty$
algebraic data (specifically the KMS temperature $\beta = 2\pi$ fixed
by Bisognano--Wichmann) that casts this shadow onto both the OTOC
Lyapunov exponent and the Krylov complexity growth rate.
\end{remark}

\paragraph{Mechanism of the shadow.}
The chain of identifications is as follows.
The ECI saturation identity (established in~\cite{ECI_v6})
gives the generalised entropy growth rate at saturation as
\begin{equation}
  \frac{1}{S_{\mathrm{gen}}} \frac{dS_{\mathrm{gen}}}{d\tau}
  \bigg|_{\mathrm{sat}}
  \;=\; \frac{2\pi T_H}{\hbar}
  \;=\; \frac{\kappa_R}{\hbar},
  \label{eq:eci-sat}
\end{equation}
where $\kappa_R = 2\pi T_H$ is the Bisognano--Wichmann surface gravity.
Under the Faulkner--Speranza identification $S_{\mathrm{gen}}
\leftrightarrow \mathcal{C}_K^{\mathrm{mod}}$ (modular flow computes
the generalised entropy~\cite{FaulknerSperanza2024}),
equation~\eqref{eq:eci-sat} becomes equation~\eqref{eq:conj-main}.
The MSS bound~\eqref{eq:mss} at saturation gives precisely the same
rate $2\pi/\beta$, completing the triangle:

\[
\underbrace{\text{Type-II}_\infty \text{ KMS at } \beta=2\pi}_{\text{algebraic input}}
\;\longrightarrow\;
\underbrace{\kappa_R = 2\pi T_H}_{\text{surface gravity}}
\;\longrightarrow\;
\underbrace{\lambda_L = 2\pi/\beta}_{\text{MSS saturation}}.
\]

We stress again: Conjecture~\ref{conj:msc} is a working conjecture.
The bridge $S_{\mathrm{gen}} \leftrightarrow \mathcal{C}_K^{\mathrm{mod}}$
is consistent with the Witten--CLPW--Caputa--Faulkner--Speranza
programme but has not been established as a theorem in the
type-II$_\infty$ crossed-product setting.
A formal derivation constitutes the main open problem, discussed in
\S\ref{sec:open}.

% ─────────────────────────────────────────────────────────────────────
\section{Verification on Tractable Models}
\label{sec:models}
% ─────────────────────────────────────────────────────────────────────

\subsection{Free scalar field on the Rindler wedge}

Consider a free massless scalar field $\hat\phi(x)$ in $d$-dimensional
Minkowski space, with the algebra of observables restricted to the
right Rindler wedge $\mathcal{W} = \{x^1 > |x^0|\}$.
By the Bisognano--Wichmann theorem, the vacuum state restricted to
$\mathcal{W}$ is KMS at $\beta_{\mathrm{BW}} = 2\pi$ with respect to
the boost Hamiltonian
\begin{equation}
  K_{\mathrm{boost}} = 2\pi \int_{\mathcal{W}} T^{00}(x)\, x^1\, d^{d-1}x,
\end{equation}
and the modular Hamiltonian is $K_\phi = 2\pi K_{\mathrm{boost}}$.

The modular flow $\sigma_s(\hat\phi(x)) = \hat\phi(\Lambda_s x)$,
where $\Lambda_s$ is the boost with rapidity $s$, generates operator
evolution in modular time.
The Krylov complexity of $\hat\phi(x)$ under this modular evolution
can be computed via the Lanczos algorithm applied to the inner product
$(\mathcal{O}, \mathcal{O}') = \langle \mathcal{O}^\dagger
\mathcal{O}'\rangle_\phi$ defined by the KMS state $\phi$.
The boost-KMS autocorrelation function is
\begin{equation}
  C(s) = \langle \hat\phi(x)\,\sigma_s(\hat\phi(x))\rangle_\phi
       = \frac{1}{(2\sinh(\pi s))^{d-2}}\cdot(\text{spatial factors}),
\end{equation}
which decays as $e^{-\pi |s|}$ at large $|s|$ in the $d=2$ case.
The Lanczos coefficients $b_n$ grow linearly as $b_n \sim n\cdot
(\pi/\beta_{\mathrm{BW}}) = n/2$ (using universal results of
Parker et al.~\cite{Parker2019}), giving an exponential Krylov
complexity growth with
\begin{equation}
  \lambda_K^{\mathrm{Rindler}} = 2b_1 = \frac{2\pi}{\beta_{\mathrm{BW}}} = 1.
  \label{eq:rindler-krylov}
\end{equation}
This is consistent with Conjecture~\ref{conj:msc}: the type-II$_\infty$
algebra of the Rindler wedge, with its canonical $\beta = 2\pi$, yields
$\lambda_K = 1$ saturating the MSS bound.

\subsection{SYK model at large $N$}

The Sachdev--Ye--Kitaev (SYK) model of $N$ Majorana fermions with
$q$-body all-to-all random interactions is a canonical maximally chaotic
system~\cite{SYK_Kitaev}. At temperature $\beta^{-1}$, the OTOC
Lyapunov exponent saturates $\lambda_L^{\mathrm{OTOC}} = 2\pi/\beta$.

For the Krylov complexity, Parker et al.~\cite{Parker2019} computed
the SYK Lanczos coefficients numerically and confirmed
$\lambda_K = 2\pi/\beta$, saturating the MSS bound.
From the algebraic perspective, the large-$N$ SYK algebra is type-III$_1$
at finite $N$ but approaches a type-II$_\infty$ structure in the
strict $N \to \infty$ limit, where the crossed-product
construction~\eqref{eq:crossed} becomes exact.

Heller, Papalini, and Schuhmann~\cite{HellerPapalini2024} establish,
via DSSYK / sine-dilaton gravity, that Krylov spread complexity
matches the complexity-equals-volume proposal at finite temperatures,
providing a quantitative holographic check.
Their results, consistent with SYK at large $N$ saturating the MSS bound, give
\begin{equation}
  \lambda_K^{\mathrm{SYK}} = \frac{2\pi}{\beta},
  \label{eq:syk-krylov}
\end{equation}
consistent with Conjecture~\ref{conj:msc}.

\begin{remark}
At finite $N$, the SYK algebra is type-III$_1$ and does not admit a
semifinite trace.
As we discuss in \S\ref{sec:implications}, Conjecture~\ref{conj:msc}
predicts that the sharpness of the MSS saturation should improve
monotonically as $N \to \infty$, tracking the convergence of the
algebra from type-III to type-II$_\infty$.
This gives an algebra-type-sensitive finite-$N$ correction
$\delta\lambda_K \sim N^{-1}$ (schematically), a quantitative
prediction that could in principle be tested numerically.
\end{remark}

% ─────────────────────────────────────────────────────────────────────
\section{Implications}
\label{sec:implications}
% ─────────────────────────────────────────────────────────────────────

The Modular Shadow Conjecture, if correct, has several structural
consequences.

\paragraph{Algebra-type classification of MSS saturation.}
The conjecture implies that MSS bound saturation is an
\emph{algebra-type phenomenon}:
\begin{itemize}
\item \textbf{Type-II$_\infty$ (crossed-product of wedge algebras at
  large $N$):} the semifinite trace and the KMS condition at $\beta=2\pi$
  together force $\lambda_K = 2\pi/\beta$, saturating the MSS bound.
\item \textbf{Type-III$_1$ (wedge algebras at finite $N$):}
  no semifinite trace; the modular flow is ergodic and the Krylov
  complexity may grow faster or fluctuate more. Saturation is not
  guaranteed.
\item \textbf{Type-I or type-II$_1$ (finite-dimensional systems or
  compact groups):} the Lyapunov exponent is bounded by the operator
  norm of the Hamiltonian; MSS saturation requires additional conditions
  (fast-scrambling dynamics) and is not a consequence of the algebraic
  structure alone.
\end{itemize}
This trichotomy gives the modular shadow conjecture predictive force:
it singles out type-II$_\infty$ algebras as the unique structural home
of MSS saturation.

\paragraph{Large-$N$ limit and algebra convergence.}
A precise form of the above is: as $N \to \infty$ in SYK, the
operator algebra converges (in the weak-operator topology) from
type-III$_1$ to type-II$_\infty$, and simultaneously $\lambda_K$
converges to $2\pi/\beta$.
The conjecture predicts that these limits are not coincidental but
represent the same structural transition.
A numerical test at $N = 10, 20, 50, 100$ in SYK would provide
evidence for or against this prediction.

\paragraph{Generalised entropy and the Faulkner--Speranza bridge.}
Faulkner and Speranza~\cite{FaulknerSperanza2024} showed that the
modular flow on the type-II$_\infty$ algebra computes the generalised
entropy $S_{\mathrm{gen}} = S_{\mathrm{BH}} + S_{\mathrm{bulk}}$.
The modular shadow conjecture provides the converse direction:
if $S_{\mathrm{gen}}$ grows with a Lyapunov exponent $2\pi/\beta$,
then the underlying algebra must be type-II$_\infty$ with the canonical
KMS temperature.
This makes $\lambda_K$ a \emph{diagnostic} of the algebraic type of
the observable algebra, in principle measurable from the entropy growth
rate alone.

\paragraph{Connes--Stoermer entropy and the type-II$_1$ reduction.}
The type-II$_\infty$ algebra can be cut down to a type-II$_1$
factor by a density matrix conditioning (passing to a finite-volume
subsystem).
In the type-II$_1$ factor, the Connes--Stoermer theorem ensures that
the one-parameter modular automorphism group preserves the unique
normalised trace~\cite{ConnesRovelli1994}.
This suggests a dual formulation of the conjecture: the Lyapunov
exponent is controlled by the modular automorphism group's action on
the Connes--Stoermer trace, and the value $2\pi/\beta$ is the unique
fixed point of this action compatible with trace-preservation.

% ─────────────────────────────────────────────────────────────────────
\section{Open Questions}
\label{sec:open}
% ─────────────────────────────────────────────────────────────────────

\begin{enumerate}
\item \textbf{Rigorous proof in Rindler.}
  The heuristic argument of \S\ref{sec:models} relies on the linear
  growth of Lanczos coefficients $b_n \sim n/2$.
  A rigorous proof would require establishing this asymptotics directly
  from the modular Hamiltonian spectrum of the free-field type-II$_\infty$
  algebra, using the framework of~\cite{HellerPapalini2024}.

\item \textbf{Bianchi vacuum and spatially homogeneous cosmologies.}
  The Bianchi~IX (mixmaster) universe exhibits chaotic dynamics.
  The modular Hamiltonian of the associated wedge algebra has been
  partially studied in the ECI programme~\cite{ECI_v6}.
  The eigenvalue-crossing gap for the Hadamard SLE on the BKL attractor
  has been closed in~\cite{BIX_lemma_A1}.
  Establishing whether the Bianchi algebra is type-II$_\infty$ and
  whether its Krylov exponent saturates~\eqref{eq:conj-main}
  would constitute a strong test of the conjecture in a cosmological
  setting.

\item \textbf{FRW and de Sitter.}
  The static patch of de Sitter space has $\beta_{\mathrm{dS}} =
  2\pi/H$ (Gibbons--Hawking temperature) and a wedge algebra conjectured
  to be type-II$_\infty$ by Witten~\cite{Witten2022}.
  Establishing whether the Krylov exponent of the de Sitter wedge
  algebra saturates $\lambda_K = H$ (the Hubble rate) would connect
  the conjecture to cosmic censorship and the dS/CFT correspondence.

\item \textbf{AdS--Schwarzschild.}
  The BTZ black hole at $\beta = 2\pi \ell_{\mathrm{AdS}}$ is the
  holographic paradigm for MSS saturation.
  A proof of Conjecture~\ref{conj:msc} in this setting would directly
  link the algebraic structure of the bulk algebra to the boundary
  conformal field theory's maximal chaos.

\item \textbf{Connection to MSS via Connes--Stoermer.}
  The original MSS proof uses analyticity of the OTOC in the complex
  time plane.
  A purely algebraic proof would bypass the large-$N$ assumption and
  derive the bound from the Tomita--Takesaki modular theory for
  type-II$_\infty$ algebras.
  The Connes--Stoermer inequality~\cite{ConnesRovelli1994} constrains
  the relative entropy between KMS states; a version of the MSS bound
  may follow from this inequality applied to modular perturbations.
  A derivation route via the Araki relative modular operator in the
  type-II$_\infty$ crossed-product setting is pursued in~\cite{EREPR_Araki}.
\end{enumerate}

% ─────────────────────────────────────────────────────────────────────
\section*{Acknowledgments}
% ─────────────────────────────────────────────────────────────────────

This note was extracted and elaborated from the ECI framework paper
(v6.0.44, Zenodo record 20021358, 2026-05-04), where the modular
shadow identification appears in \S\ref{sec:mss-shadow}.
The conjecture status was confirmed by a multi-agent cross-check
(2026-04-30).
We thank [collaborators] for discussions.

% ─────────────────────────────────────────────────────────────────────
\begin{thebibliography}{99}
% ─────────────────────────────────────────────────────────────────────

\bibitem{MSS2016}
J.~Maldacena, S.~H.~Shenker, and D.~Stanford,
``A bound on chaos,''
\textit{JHEP} \textbf{08} (2016) 106,
\href{https://doi.org/10.1007/JHEP08(2016)106}{doi:10.1007/JHEP08(2016)106},
\href{https://arxiv.org/abs/1503.01409}{arXiv:1503.01409 [hep-th]}.

\bibitem{Parker2019}
D.~E.~Parker, X.~Cao, A.~Avdoshkin, T.~Scaffidi, and E.~Altman,
``A Universal Operator Growth Hypothesis,''
\textit{Phys. Rev. X} \textbf{9} (2019) 041017,
\href{https://doi.org/10.1103/PhysRevX.9.041017}{doi:10.1103/PhysRevX.9.041017},
\href{https://arxiv.org/abs/1812.08657}{arXiv:1812.08657 [cond-mat.str-el]}.

\bibitem{Caputa2024}
P.~Caputa, J.~M.~Mag\'{a}n, D.~Patramanis, and E.~Tonni,
``Krylov complexity of modular Hamiltonian evolution,''
\textit{Phys. Rev. D} \textbf{109} (2024) 086004,
\href{https://doi.org/10.1103/PhysRevD.109.086004}{doi:10.1103/PhysRevD.109.086004},
\href{https://arxiv.org/abs/2306.14732}{arXiv:2306.14732 [hep-th]}.

\bibitem{Witten2022}
E.~Witten,
``Gravity and the crossed product,''
\textit{JHEP} \textbf{10} (2022) 008,
\href{https://doi.org/10.1007/JHEP10(2022)008}{doi:10.1007/JHEP10(2022)008},
\href{https://arxiv.org/abs/2112.12828}{arXiv:2112.12828 [hep-th]}.

\bibitem{CLPW2023}
V.~Chandrasekaran, R.~Longo, G.~Penington, and E.~Witten,
``An algebra of observables for de~Sitter space,''
\textit{JHEP} \textbf{02} (2023) 082,
\href{https://doi.org/10.1007/JHEP02(2023)082}{doi:10.1007/JHEP02(2023)082},
\href{https://arxiv.org/abs/2206.10780}{arXiv:2206.10780 [hep-th]}.

\bibitem{FaulknerSperanza2024}
T.~Faulkner and A.~J.~Speranza,
``Gravitational algebras and the generalized second law,''
\textit{JHEP} \textbf{11} (2024) 099,
\href{https://doi.org/10.1007/JHEP11(2024)099}{doi:10.1007/JHEP11(2024)099},
\href{https://arxiv.org/abs/2405.00847}{arXiv:2405.00847 [hep-th]}.

\bibitem{DEHK2025a}
J.~De~Vuyst, S.~Eccles, P.~A.~H\"{o}hn, and J.~Kirklin,
``Crossed products and quantum reference frames: on the observer-dependence
of gravitational entropy,''
\textit{JHEP} \textbf{07} (2025) 063,
\href{https://doi.org/10.1007/JHEP07(2025)063}{doi:10.1007/JHEP07(2025)063},
\href{https://arxiv.org/abs/2412.15502}{arXiv:2412.15502 [hep-th]}.

\bibitem{DEHK2025b}
J.~De~Vuyst, S.~Eccles, P.~A.~H\"{o}hn, and J.~Kirklin,
``Gravitational entropy is observer-dependent,''
\textit{JHEP} \textbf{07} (2025) 146,
\href{https://doi.org/10.1007/JHEP07(2025)146}{doi:10.1007/JHEP07(2025)146},
\href{https://arxiv.org/abs/2405.00114}{arXiv:2405.00114 [hep-th]}.

\bibitem{HellerPapalini2024}
M.~P.~Heller, A.~Papalini, and T.~Schuhmann,
``Krylov spread complexity $\leftrightarrow$ complexity = volume in DSSYK / sine-dilaton gravity at finite T,''
\textit{Phys. Rev. Lett.} \textbf{135} (2025) 151602,
\href{https://doi.org/10.1103/PhysRevLett.135.151602}{doi:10.1103/PhysRevLett.135.151602},
\href{https://arxiv.org/abs/2412.17785}{arXiv:2412.17785 [hep-th]}.
%% [FIXED 2026-05-04: HPS paper is about DSSYK/sine-dilaton C=V matching at finite T,
%%  NOT about ``Krylov complexity of modular Hamiltonian evolution''. That description
%%  belongs to Caputa-Magan-Patramanis-Tonni arXiv:2306.14732 (cited as \cite{Caputa2024}).]

\bibitem{BisognanoWichmann1976}
J.~J.~Bisognano and E.~H.~Wichmann,
``On the duality condition for a Hermitian scalar field,''
\textit{J. Math. Phys.} \textbf{17} (1976) 303.

\bibitem{ConnesRovelli1994}
A.~Connes and C.~Rovelli,
``Von Neumann algebra automorphisms and time-thermodynamics relation in
generally covariant quantum theories,''
\textit{Class. Quantum Grav.} \textbf{11} (1994) 2899--2917,
\href{https://doi.org/10.1088/0264-9381/11/12/007}{doi:10.1088/0264-9381/11/12/007}.

\bibitem{SYK_Kitaev}
A.~Kitaev, ``A simple model of quantum holography,''
Talks at KITP, 7 and 23 April 2015,
\url{http://online.kitp.ucsb.edu/online/entangled15/kitaev/}.

\bibitem{ECI_v6}
[ECI author(s)],
``Entanglement--Complexity--Information: a phenomenological architecture
for quantum gravity phenomenology (v6.0.44),''
Zenodo record 20021358, 2026-05-04.

\bibitem{CardyRho}
[ECI author(s)],
``Universal Analog-Hawking Saturation Ratio $\rho = c/12$
for Unitary Diagonal Minimal-Invariant-Partition CFTs,''
arXiv:[ECI v6.2 Cardy paper, in preparation, 2026].
%% Note: ``\rho'' in that paper is the analog-Hawking saturation ratio, not a density matrix.

\bibitem{BIX_lemma_A1}
K.~Remondi\`{e}re,
``Appendix A: Closure of the Eigenvalue-Crossing Gap for the Hadamard SLE on Bianchi~IX,''
arXiv:[ECI v6.2 Bianchi~IX, in preparation, 2026].

\bibitem{EREPR_Araki}
[ECI author(s)],
``Araki Relative Modular Operator in Type-II$_\infty$: Week 2 Scoping Memo
(ER=EPR Reopen, A13-2),''
arXiv:[ECI v6.2 ER=EPR route, in preparation, 2026].

\end{thebibliography}

\end{document}
